\begin{textsample}{POS Dim 1 – global\_south\_2024\_10 – Score 32.00 – global\_south\_2024\_10/116710388.txt}  \label{ex:f1_pos_166}
The prolonged conflict has also overwhelmed local \textbf{emergency} responders, who are struggling to meet the urgent needs of the \textbf{civilian} \textbf{population}.
A Palestinian man walks past the rubble after Israeli forces withdrew from the \textbf{area} around Kamal Adwan \textbf{hospital}, amid the ongoing Israel-Hamas conflict, in Jabalia, in the \textbf{northern} Gaza Strip October 26, 2024. ( REUTERS Photo ) On Monday, Israeli \textbf{tanks} launched a significant incursion into \textbf{northern} Gaza, penetrating two towns and a historic \textbf{refugee} \textbf{camp}, which resulted in approximately 100,000 \textbf{civilians} being trapped, according to reports from the Palestinian \textbf{emergency} service.
This military action is part of a broader Israeli \textbf{operation} aimed at dismantling \textbf{regrouping} Hamas militants in the region, further escalating the already dire \textbf{humanitarian} situation in Gaza.
The Israeli military reported that soldiers successfully apprehended around 100 suspected Hamas militants during a raid on the Kamal Adwan \textbf{hospital} located within the Jabalia \textbf{refugee} \textbf{camp}.
However, both Hamas and \textbf{medical} personnel at the \textbf{hospital} have denied the presence of any militants in the \textbf{facility}.
This contradiction underscores the complexities and contentious narratives surrounding the conflict @ @ @ @ @ @ @ @ @ @ at least 19 individuals were killed due to Israeli airstrikes and \textbf{bombardment} on Monday alone, with 13 of these casualties occurring in the \textbf{northern} parts of the \textbf{territory}.
The Palestinian Civil \textbf{Emergency} Service has reported that the ongoing conflict has left around 100,000 residents marooned in Jabalia, Beit Lahiya, and Beit Hanoun, facing \textbf{severe} \textbf{shortages} of \textbf{medical} \textbf{supplies} and \textbf{food}.
This alarming statistic has yet to be independently verified by Reuters, but it highlights the critical \textbf{humanitarian} crisis unfolding in these \textbf{areas}.
Advertisement \textbf{Emergency} services reported that their \textbf{operations} have come to a standstill due to the three-week-long Israeli \textbf{offensive}, which has focused on regions previously thought to have eliminated significant Hamas combat forces.
The prolonged conflict has overwhelmed local \textbf{emergency} responders, who are struggling to meet the urgent needs of the \textbf{civilian} \textbf{population}.
As military actions \textbf{intensified}, diplomatic discussions aimed at brokering a ceasefire resumed on Sunday, led by the United States, Egypt, and Qatar.
After multiple unsuccessful attempts, Egyptian President Abdel Fattah el-Sisi proposed a preliminary two-day @ @ @ @ @ @ @ @ @ @ held by Hamas for a number of Palestinian prisoners.
Following this initial agreement, discussions could take place within ten days to negotiate a more lasting ceasefire.
However, there has been no public comment from either Israel or Hamas regarding these proposals, as both parties continue to adhere to fundamentally opposing conditions for ending the conflict.
This impasse highlights the difficulties in achieving a resolution, with both sides remaining entrenched in their positions.
The ongoing conflict in Gaza has heightened tensions across the Middle East, raising fears of a more widespread regional instability.
Israeli military \textbf{operations} have extended into \textbf{southern} Lebanon in response to Hezbollah ' s rocket attacks on \textbf{northern} Israel, aimed at supporting Hamas.
This situation has not only exacerbated the \textbf{humanitarian} crisis in Gazabut has also triggered rare direct confrontations between Israel and Iran.
Over the weekend, Israeli warplanes conducted strikes on missile sites in Iran, retaliating for an Iranian missile attack on Israel that occurred on October 1.
Iran ' s Foreign \textbf{Ministry} responded by declaring that Tehran would"use @ @ @ @ @ @ @ @ @ @ actions.
Advertisement The \textbf{humanitarian} crisis in \textbf{northern} Gaza is further compounded by the dire state of local \textbf{healthcare} \textbf{facilities}.
Reports indicate that \textbf{hospitals} in the region are struggling to operate under extreme conditions.
Despite the Israeli military ' s insistence on evacuating \textbf{hospitals}, officials at these \textbf{facilities} have refused such orders.
At least two \textbf{hospitals} have been damaged due to Israeli fire, resulting in critical \textbf{shortages} of \textbf{medical} \textbf{supplies}, \textbf{food}, and \textbf{fuel}.
In a tragic reflection of the situation, the Gaza \textbf{health} \textbf{ministry} reported that in the past week, at least one doctor, a nurse, and two \textbf{child} \textbf{patients} have died in \textbf{hospitals} due to a lack of \textbf{medical} treatment.
Currently, only one \textbf{medical} staff member -- a pediatrician -- remains at Kamal Adwan \textbf{Hospital} after Israeli forces detained and expelled the remaining \textbf{healthcare} \textbf{workers}.
This highlights the \textbf{severe} impact of military \textbf{operations} on \textbf{medical} care and the urgent need for \textbf{humanitarian} assistance.
Residents in \textbf{northern} Gaza have reported increasingly desperate conditions as Israeli forces \textbf{besiege} schools and \textbf{shelters} where @ @ @ @ @ @ @ @ @ @ soldiers compelling families to evacuate these makeshift \textbf{shelters} while rounding up men and directing women and \textbf{children} towards Gaza \textbf{City} and \textbf{southern} \textbf{areas} of the \textbf{territory}.
This tactic has led to a pervasive sense of fear and uncertainty among the \textbf{civilian} \textbf{population}, who find themselves caught in the crossfire of an intensifying military campaign.

% matched lemmas: area, besiege, bombardment, camp, child, city, civilian, emergency, facility, food, fuel, health, healthcare, hospital, humanitarian, intensified, medical, ministry, northern, offensive, operation, patient, population, refugee, regroup, severe, shelter, shortage, southern, supply, tank, territory, worker
\end{textsample}
